% !TeX program = xelatex

\documentclass[UTF8,a4paper,zihao=-4,fontset=windows]{ctexart}
% Shared renderer entrypoint for generated lessonplan.tex files.
% Generated documents should load ctexart, then input this file.

\usepackage{hepingjie_lessonplan}
% Hepingjie lesson-plan overlay coordinates.
% Coordinate origin: top-left corner of an A4 portrait page.
% Units: mm. These coordinates are frozen for templates/lessonplan/hepingjie_blank.pdf.

\newcommand{\LPBackgroundPdf}{templates/lessonplan/hepingjie_blank.pdf}

% Page 1 fixed fields.
\newcommand{\LPDateYearX}{65mm}
\newcommand{\LPDateMonthX}{89mm}
\newcommand{\LPDateDayX}{110mm}
\newcommand{\LPDateY}{38mm}
\newcommand{\LPRecordPageX}{170mm}
\newcommand{\LPRecordPageY}{38mm}

\newcommand{\LPTopicX}{32mm}
\newcommand{\LPTopicY}{49mm}
\newcommand{\LPTopicW}{108mm}
\newcommand{\LPTypeX}{154mm}
\newcommand{\LPTypeY}{49mm}
\newcommand{\LPTypeW}{40mm}

\newcommand{\LPUnitTotalX}{67mm}
\newcommand{\LPUnitTotalY}{61mm}
\newcommand{\LPTopicTotalX}{111mm}
\newcommand{\LPTopicTotalY}{61mm}
\newcommand{\LPCurrentPeriodX}{169mm}
\newcommand{\LPCurrentPeriodY}{61mm}

\newcommand{\LPBodyX}{33mm}
\newcommand{\LPObjectiveY}{74mm}
\newcommand{\LPObjectiveH}{39mm}
\newcommand{\LPKeyPointY}{119mm}
\newcommand{\LPKeyPointH}{15mm}
\newcommand{\LPDifficultyY}{137mm}
\newcommand{\LPDifficultyH}{15mm}
\newcommand{\LPMethodY}{156mm}
\newcommand{\LPMethodH}{15mm}
\newcommand{\LPMediaY}{174mm}
\newcommand{\LPMediaH}{15mm}
\newcommand{\LPBlackboardY}{196mm}
\newcommand{\LPBlackboardH}{73mm}
\newcommand{\LPBodyW}{164mm}

% Teaching process boxes on pages 2-4.
\newcommand{\LPProcessX}{33mm}
\newcommand{\LPProcessY}{31mm}
\newcommand{\LPProcessW}{163mm}
\newcommand{\LPProcessFullH}{249mm}
\newcommand{\LPProcessPageFourH}{160mm}


\usepackage{tikz}

\begin{document}

\LPBackgroundPage{1}{
  \LPCenterBox{\LPDateYearX}{\LPDateY}{12mm}{5mm}{}
  \LPCenterBox{\LPDateMonthX}{\LPDateY}{10mm}{5mm}{}
  \LPCenterBox{\LPDateDayX}{\LPDateY}{10mm}{5mm}{}
  \LPCenterBox{\LPRecordPageX}{\LPRecordPageY}{10mm}{5mm}{1}
  \LPTextBox{\LPTopicX}{\LPTopicY}{\LPTopicW}{8mm}{移项解一元一次方程}
  \LPTextBox{\LPTypeX}{\LPTypeY}{\LPTypeW}{8mm}{新授课}
  \LPCenterBox{\LPUnitTotalX}{\LPUnitTotalY}{12mm}{5mm}{8}
  \LPCenterBox{\LPTopicTotalX}{\LPTopicTotalY}{12mm}{5mm}{1}
  \LPCenterBox{\LPCurrentPeriodX}{\LPCurrentPeriodY}{12mm}{5mm}{1}
  \LPTextBox{\LPBodyX}{\LPObjectiveY}{\LPBodyW}{\LPObjectiveH}{\LPItems{\item 能说明移项来源于等式性质，而不是只背口诀。
\item 能按移项、合并同类项、系数化为 1 的步骤解方程。
\item 能处理移项变号和负系数，并完成代入检验。
\item 能用一元一次方程解决简单实际问题。
}}
  \LPTextBox{\LPBodyX}{\LPKeyPointY}{\LPBodyW}{\LPKeyPointH}{移项法则及其在解一元一次方程中的规范应用。}
  \LPTextBox{\LPBodyX}{\LPDifficultyY}{\LPBodyW}{\LPDifficultyH}{理解移项必须变号的依据，并在负系数和实际问题中保持符号正确。}
  \LPTextBox{\LPBodyX}{\LPMethodY}{\LPBodyW}{\LPMethodH}{复习迁移; 对比推理; 板演示范; 错解分析; 反馈检测}
  \LPTextBox{\LPBodyX}{\LPMediaY}{\LPBodyW}{\LPMediaH}{PPT; 黑板; 流程卡; 反馈卡}
  \LPTextBox{\LPBodyX}{\LPBlackboardY}{\LPBodyW}{\LPBlackboardH}{\begin{enumerate}[leftmargin=1.4em,itemsep=0.4mm,topsep=0pt]
\item 移项依据：等式两边同加或同减同一个式子。
\item 移项规则：项从一边到另一边，符号改变。
\item 步骤：移项，合并同类项，系数化为 1，检验。
\item 例1：$3x+7=32-2x$。
\item 例2：$x-3=1.5x+1$。
\item 易错：移项不变号，负系数丢负号。
\end{enumerate}}
}

\LPBackgroundPage{2}{
  \LPProcessBox{\LPProcessX}{\LPProcessY}{\LPProcessW}{\LPProcessFullH}{\LPProcessRow{复习与问题引入}{{\heiti 导入题}\par 问题1：解 $x+5=12$，每一步依据是什么？预设：两边同时减 5。问题2：方程 $3x+7=32-2x$ 中，未知数项分在两边，怎样把它们集中？教师不直接给口诀，而是要求先用等式性质完整变形。}{\LPStudentItems{\item 解热身方程并说依据。
\item 观察新方程中未知数项分居两边。
\item 提出需要把同类项集中。
}}{5}
\LPProcessRow{移项依据与基础例题}{{\heiti 例题1 从等式性质推出移项}\par 题目：解 $3x+7=32-2x$。关键问：要消去右边的 $-2x$，两边同时加什么？规范解：\[\begin{aligned}3x+7&=32-2x,\\3x+2x&=32-7,\\5x&=25,\\x&=5.\end{aligned}\] 检验：左边 $3\times5+7=22$，右边 $32-2\times5=22$。即时练习：$4x-3=2x+9$，答：$x=6$。机动追问：若把 $7$ 移到右边仍写 $+7$，等式是否还成立？答：不成立。}{\LPStudentItems{\item 圈出要移动的项。
\item 先写完整等式性质变形，再写移项简写。
\item 完成即时练习并检验。
}}{10}
\LPProcessRow{变式例题与负系数}{{\heiti 例题2 负系数处理}\par 题目：解 $x-3=1.5x+1$。教师问：把 $1.5x$ 移到左边后左边是多少？规范解：\[\begin{aligned}x-3&=1.5x+1,\\x-1.5x&=1+3,\\-0.5x&=4,\\x&=-8.\end{aligned}\] 强调：负系数化为 1 时两边同除以 $-0.5$，负号不能丢。即时练习：$2x+1=5x-8$，答：$x=3$。备用追问：若得到 $-2x=6$，解是 $3$ 还是 $-3$？答：$-3$。}{\LPStudentItems{\item 独立完成前两步。
\item 说明为什么 $-0.5x=4$ 的解是 $-8$。
\item 用代入法验证。
}}{8}}
}

\LPBackgroundPage{2}{
  \LPProcessBox{\LPProcessX}{\LPProcessY}{\LPProcessW}{\LPProcessFullH}{\LPProcessRow{易错辨析与纠错}{{\heiti 例题3 易错辨析}\par 错解：由 $2x-4=12$ 得 $2x=8$，所以 $x=4$。教师问：$-4$ 从左边到右边应变成什么？规范纠错：$2x=12+4$，$2x=16$，$x=8$。检验：$2\times8-4=12$。即时练习：判断 $5x+6=2x$ 移项为 $5x-2x=-6$ 是否正确，答：正确。}{\LPStudentItems{\item 判断错解的第一处错误。
\item 把被移动的项连同符号一起圈出。
\item 订正后代入原方程。
}}{7}
\LPProcessRow{应用例题}{{\heiti 例题4 应用}\par 题目：甲、乙两组共整理 55 本图书，甲组比乙组的 2 倍少 5 本。设乙组整理 $x$ 本，求乙组本数。关键问：甲组怎样表示？规范解：甲组 $2x-5$，方程 $x+2x-5=55$，$3x=60$，$x=20$。答：乙组整理 20 本。即时练习：一数的 3 倍比这个数多 18，求这个数，答：9。}{\LPStudentItems{\item 列出相等关系。
\item 写设句、方程、解和答语。
\item 互查每次移项是否变号。
}}{6}
\LPProcessRow{随堂练习与补救}{{\heiti 练习组与补救}\par 必做：1. $5x+2=3x+10$。2. $4-x=2x-5$。3. 判断 $6-2x=10$ 移项得 $-2x=4$ 是否正确。答案：$x=4$；$x=3$；正确。备用练习1：解 $3-2x=5x+17$，解：$-2x-5x=17-3$，$-7x=14$，$x=-2$。备用练习2：判断 $4x+7=2x-5$ 得 $4x-2x=-5-7$，$2x=-12$，$x=-6$ 是否正确。答案：正确，两次移项均改变符号。巡视关注：$3$ 移到右边变 $-3$，$5x$ 移到左边变 $-5x$，负系数化为 1 不丢负号。备用题只在学生完成较快时使用，不另计固定时间。}{\LPStudentItems{\item 独立完成三题。
\item 同桌互判并标出变号位置。
\item 错误学生回写等式性质依据。
}}{4}}
}

\LPBackgroundPage{4}{
  \LPProcessBox{\LPProcessX}{\LPProcessY}{\LPProcessW}{\LPProcessPageFourH}{\LPProcessRow{反馈测试、小结与作业}{{\heiti 课堂反馈测试 5分}\par 共 10 分。1. 移项为什么要变号？2分。2. 解 $6x+5=2x+17$，3分。3. 解 $0.5x-2=x+1$，3分。4. 判断 $3x-7=8$ 移项得 $3x=1$ 是否正确并改正，2分。答案：依据等式性质；$x=3$；$x=-6$；错误，应为 $3x=15,x=5$。同桌互判后举手统计，90\%以上做提高题，70\%以下重做“完整等式性质写法”。\par {\heiti 分层作业}\par A组：1. $2x+3=11$。2. $5x-4=2x+8$。3. $7-x=3x-5$。4. 判断 $x-6=9$ 移项为 $x=3$ 的错因。5. 检验 $x=4$ 是否为 $3x+1=13$ 的解。B组：1. $0.2x+1=0.5x-5$。2. $\frac12x-3=x+2$。3. 用方程解决：两个连续整数和为 37，求这两个数。C组选做：写一段话解释移项与等式性质的关系。答案：A1 $4$，A2 $4$，A3 $3$，A4 应变为 $+6$，A5 是；B1 $20$，B2 $-10$，B3 $18,19$。}{\LPStudentItems{\item 独立完成反馈卡。
\item 同桌互判后举牌反馈。
\item 说出一句话：移项是等式性质的简写。
\item 记录分层作业。
}}{5}}
}

\end{document}
